% M25, 13.09.2026: bestehende Dokumentationsbelege in den numerischen Zitierstil überführt.
% M10, 11.09.2026: Leserführung, Kapiteltrennung und belegte Reichweite.
% PipelineFullRunner.cs / .Cli.cs / .EventPump.cs; IngestionApplyExec.cs:31--41;
% GithubForwardApply.cs:69--86; GateFreeBinding.cs; SubworkflowPortSpikeTests.cs.
% B-33: 13 Checkpoints des Elternlaufs bis zur Pause, nicht 13 Ledger-Checkpoints.
% BindGateFree weist alle Ports zurück. AsAIAgent-Dokumentation ist kein BindAsExecutor-Beleg.
% Historischer Fall und Prozessgrenzen-Smoke in 7.6; keine neue Systemausführung für M10.

\section{Durabilität: Checkpoints, Wiederaufnahme und Teilworkflows}
\label{sec:realisierung-durabilitaet}

MAF sichert den Workflowfortschritt an abgeschlossenen Ausführungsrunden
(\emph{Supersteps}) in Checkpoints. Das Artefakt verwendet dafür einen
dateibasierten JSON-Speicher. Zur Wiederaufnahme wird der Graph erneut
aufgebaut und der gespeicherte Ausführungszustand geladen. Nach der
Frameworkdokumentation müssen dabei Graphstruktur und
Executor-Identitäten zum gespeicherten Stand passen; für neu erzeugte
ChatClient-basierte Agenten müssen auch Kennung und gegebenenfalls
Name erhalten bleiben \cite{msmaf2026checkpoints}.

Abbildung~\ref{fig:checkpoint-core-trennung} zeigt den in
Abschnitt~\ref{sec:realisierung-human-gates} beschriebenen Pausenzyklus.
Checkpoint und Lauf-Zeiger ermöglichen die Fortsetzung nach einem
Prozessende. Sie stellen dabei den Ausführungszusammenhang wieder her,
keinen früheren Core-Stand. Im Lauf mitgeführte fachliche Daten bleiben
von der maßgeblichen Speicherung des Projektzustands zu unterscheiden.

\begin{figure}[htbp]
  \centering
  \includegraphics[width=\textwidth]{figures/06/Abbildung6.5.pdf}
  \caption[Wiederaufnahme eines pausierten Workflows]{Wiederaufnahme eines pausierten Workflows: Der Checkpoint
    sichert den Ausführungszustand. Der maßgebliche Core wird dadurch
    nicht zurückgesetzt; seine Fortschreibung erfolgt über die
    vorgesehenen Governance- und Speicheroperationen.}
  \label{fig:checkpoint-core-trennung}
\end{figure}

Dieser Mechanismus trägt den geregelten Pausenfall. Ein fachlicher
Abbruch oder technischer Fehler eröffnet dagegen nicht automatisch
denselben Fortsetzungsweg (Tabelle~\ref{tab:beendigungsformen}). Für den
beschriebenen Einzelbediener-Betrieb wird zudem angenommen, dass während
einer Pause keine anderweitige Core-Änderung erfolgt. Eine durchgängige
Aktualitätsprüfung früher erzeugter Vorschläge ist nicht realisiert;
die Annahme ist keine nachträglich geprüfte Versuchsbedingung der
historischen Läufe.

\begin{table}[htbp]
  \centering
  \small
  \begin{tabular}{p{0.20\textwidth}p{0.33\textwidth}p{0.36\textwidth}}
    \hline
    \textbf{Beendigungsform} & \textbf{Gespeichert bis dahin} & \textbf{Weg zurück} \\
    \hline
    geregelte Gate-Pause & Checkpoint, Lauf-Zeiger und Stufen-Artefakte & regulärer Wiederaufnahmebefehl \\
    fachlicher Abbruch & bereits gesicherte Checkpoints und Stufen-Artefakte & neuer Lauf; keine automatische Fortsetzung \\
    technischer Fehler & zuletzt gesicherter Checkpoint und vorhandene Artefakte & nicht systematisch untersucht; kein allgemeiner Wiederherstellungsanspruch \\
    \hline
  \end{tabular}
  \caption[Beendigungsformen und Fortsetzungswege]{Beendigungsformen und Fortsetzungswege. Die geregelte
    Gate-Pause ist der realisierte und belegte Wiederaufnahmefall.}
  \label{tab:beendigungsformen}
\end{table}

Die Fortsetzung ergänzt ein projektspezifischer
\emph{Wiederholungsschutz}. Bei erneuter Anwendung eines
Einordnungsplans werden dessen Plankennung, Anwendungsvermerk und
gespeicherter Ergebnisbericht geprüft. Wurde dieser Plan bereits
übernommen, wird der Bericht zurückgegeben, ohne den Core erneut zu
schreiben. Im GitHub-Pfad verhindert eine vorhandene Issue-Zuordnung
die erneute Anlage beziehungsweise dieselbe Verknüpfung. Aktualisierungen
und Kommentare sind von diesem Kurzschluss ausgenommen: Hier ist die
Zuordnung Voraussetzung der Handlung, kein Beleg ihrer Ausführung.
Daraus folgt keine allgemeine Garantie, dass jede Außenwirkung
genau einmal eintritt, insbesondere nicht bei einem Absturz zwischen externem
Schreiben und lokaler Protokollierung. Der Kontrolltest in
Abschnitt~\ref{sec:eval-kontrollen} prüft die erneute Anwendung
derselben Plankennung am Core-Pfad.

Für die Einbettung größerer Verarbeitungseinheiten verwendet das
Artefakt \emph{gebundene Teilworkflows}, in der Implementierung auch
Kapseln genannt. So wird der Ledger-Workflow über
\texttt{BindAsExecutor} in den Gesamtgraphen eingebunden; der geprüfte
Elterncheckpoint enthält auch den inneren Checkpointbestand. Eine
persistierte Zuordnung hält seine Teillauf-Kennung fest, damit der
erneute Graphaufbau auf denselben Laufordner zugreift. Im dokumentierten
Fall blieb dadurch der bereits abgeschlossene Ledger-Teillauf bei der
Wiederaufnahme zugeordnet. Die Fortsetzung innerhalb einer noch
laufenden Ledger-Verarbeitung wurde damit nicht gezeigt. Diesen
Demonstrationsfall und den separaten technischen Nachweis der
Wiederaufnahme in einem neuen Prozess unterscheidet
Abschnitt~\ref{sec:eval-resume}.

Die Kapselung erforderte eine weitere Strukturentscheidung. In einem
isolierten Integrationstest mit MAF~1.15.0 erreichte eine innere
Portanfrage ohne Weiterleitungskanten den Eltern-Runner nicht; die
Ausführung wartete ohne sichtbares Fehlerereignis. Mit expliziter
Weiterleitung von Anfrage und Antwort zwischen Teilworkflow und
Elternport gelang derselbe Rundweg. Die produktive Bindefunktion
\texttt{BindGateFree} lässt deshalb bewusst nur portfreie
Teilworkflows zu; sie weist \emph{jeden} enthaltenen Request Port
bereits beim Graphaufbau zurück. Entscheidungspunkte bleiben in
dieser gewählten Form auf der übergeordneten Ebene.
Verhaltenstests sichern beide untersuchten Konfigurationen ab.
Die aktuelle Dokumentation beschreibt eine transparente Weiterleitung
innerer Anfragen
\cite[Abschnitt~„Requests and Responses“]{msmaf2026subworkflows}.
Der Testbefund bleibt deshalb auf die
untersuchte Version und Verdrahtung begrenzt. Er begründet die gewählte
Struktur, kein generelles Verbot von Ports in Teilworkflows.
